% !TeX root = ../main.tex
\section{Recursive types}
\label{sec:recursive}

The inductive types of Section~\ref{sec:types} are defined by constants and
constructors that refer to the type being defined: $0 : \mathbb{N}$ and
$s : \mathbb{N} \to \mathbb{N}$. This section, which goes beyond the slides, shows how a
type system expresses that self-reference. A list of integers is either empty
or an integer followed by another list:
\[
\alpha = \unit + (\tyint \times \alpha).
\]
With the types of STLC this equation has no solution. Expanding it never ends,
$\alpha = \unit + \tyint \times (\unit + \tyint \times (\dots))$, for the same reason that
$t_f = t_f \to \mathit{Int}$ failed the occurs check in Section~\ref{sec:inference}.

\subsection{$\mu$-types}

A new type former names the solution: $\mu\alpha.\, \tau$ is the type $\alpha$ that
satisfies $\alpha = \tau$, where $\tau$ may mention $\alpha$.
\[
\mathbf{List}\ \tyint = \mu\alpha.\, \unit + (\tyint \times \alpha)
\qquad
\mathbf{Tree} = \mu\alpha.\, \tyint + (\alpha \times \tyint \times \alpha)
\qquad
\mathbb{N} = \mu\alpha.\, \unit + \alpha
\]
In the last one, $\unit$ plays the role of $0$ and the right side of the sum
the role of $s$.

\paragraph{Fold and unfold.} The values of $\mu\alpha.\, \tau$ are related to its
unfolding $\tau[\mu\alpha.\tau / \alpha]$ (the body with the whole type put back for
$\alpha$) by two operations:
\begin{mathpar}
\inferrule*[right=T-Fold]{\Gamma \vdash e : \tau[\mu\alpha.\tau / \alpha]}{\Gamma \vdash \mathbf{fold}(e) : \mu\alpha.\, \tau}
\and
\inferrule*[right=T-Unfold]{\Gamma \vdash e : \mu\alpha.\, \tau}{\Gamma \vdash \mathbf{unfold}(e) : \tau[\mu\alpha.\tau / \alpha]}
\end{mathpar}
$\mathbf{fold}$ packages one layer (what a constructor such as \textit{cons}
does) and $\mathbf{unfold}$ opens one layer (what a pattern match does).
Figure~\ref{fig:folds} shows the list $[1, 2]$ built with three folds, one per
constructor.

\begin{figure}[h]
\centering
\begin{forest}
  for tree={font=\small, s sep=4mm, l sep=4mm},
  [fold
    [{$\mathit{inr}$}
      [$\times$ [1]
        [fold
          [{$\mathit{inr}$}
            [$\times$ [2] [fold [{$\mathit{inl}$} [$()$]]]]
          ]
        ]
      ]
    ]
  ]
\end{forest}
\caption{The list $[1, 2]$ as a value of $\mu\alpha.\, \unit + (\tyint \times \alpha)$. $\mathit{inl}$ and $\mathit{inr}$ choose a side of the sum.}
\label{fig:folds}
\end{figure}

Languages with explicit fold and unfold, and a type equality that only
compares syntax, are \emph{iso-recursive}; languages where $\mu\alpha.\, \tau$ and its
unfolding are the same type are \emph{equi-recursive}, and need a more complex
equality on infinite trees. ML-family languages are iso-recursive, which fits
their inference.

\subsection{Algebraic data types}

The data types of Haskell are recursive types where the constructors are the
folds and pattern matching is the unfold:
\coderange{Haskell}{haskell/RecursiveTypes.hs}{types}
\coderange{Haskell}{haskell/RecursiveTypes.hs}{functions}
\texttt{Cons} takes an element and a list (one unfolding of the type) and
returns a list; the equation \texttt{len (Cons \_ xs)} opens one layer.
Recursive types do not reopen what the occurs check closed: ordinary type
variables still cannot refer to themselves, and self-reference is allowed
only where a $\mu$ (a \texttt{data} declaration) introduces it.

\subsection{Exercises}

\begin{exercise}
Unfold $\mathbf{List}\ \tyint$ one level.
\end{exercise}

\paragraph{Solution.} $\unit + (\tyint \times \mu\alpha.\, \unit + (\tyint \times \alpha))$, that is,
$\unit + (\tyint \times \mathbf{List}\ \tyint)$: empty, or an integer and a list.

\begin{exercise}
Give a $\mu$-type for nonempty lists of booleans.
\end{exercise}

\paragraph{Solution.} $\mu\alpha.\, \tybool + (\tybool \times \alpha)$: the base case is a single
boolean instead of $\unit$.

\begin{exercise}
Why does Rust reject \texttt{enum List<T> \{ Nil, Cons(T, List<T>) \}} and
accept \texttt{Cons(T, Box<List<T>>)}?
\end{exercise}

\paragraph{Solution.} Rust stores the fields of an \texttt{enum} inline, so a
\texttt{List} would contain a \texttt{List} and have infinite size.
\texttt{Box} stores a pointer, of fixed size, to a list on the heap. Haskell
and ML represent recursive values with pointers by default, which is why
they need no annotation.
